% Standalone problem document, generated from the adjacent README.md.
% Compile with XeLaTeX. Mathematical content is maintained in Markdown.
\documentclass[11pt,a4paper]{article}
\usepackage[fontsize=10.5pt]{scrextend}
\usepackage[margin=22mm,headheight=15pt,headsep=9mm,footskip=11mm]{geometry}
\usepackage{fontspec}
\setmainfont{texgyrepagella}[Extension=.otf,UprightFont=*-regular,BoldFont=*-bold,ItalicFont=*-italic,BoldItalicFont=*-bolditalic]
\setsansfont{texgyreheros}[Extension=.otf,UprightFont=*-regular,BoldFont=*-bold,ItalicFont=*-italic,BoldItalicFont=*-bolditalic]
\setmonofont{texgyrecursor}[Extension=.otf,UprightFont=*-regular,BoldFont=*-bold,ItalicFont=*-italic,BoldItalicFont=*-bolditalic,Scale=MatchLowercase]
\usepackage{amsmath,amssymb}
\usepackage{unicode-math}
\setmathfont{texgyrepagella-math.otf}
\usepackage{xcolor,microtype,enumitem,needspace,fancyhdr}
\usepackage{bookmark}
\definecolor{ink}{HTML}{17344B}
\definecolor{accent}{HTML}{197D80}
\definecolor{muted}{HTML}{596773}
\hypersetup{colorlinks=true,linkcolor=accent,urlcolor=accent,pdftitle={SP-08 - Rank-two
maximizers of spectral spread on an entry
interval},pdfauthor={Open Problems in Numerical Linear Algebra}}
\urlstyle{same}
\setlength{\parindent}{0pt}
\setlength{\parskip}{5pt plus 1pt minus 1pt}
\linespread{1.04}
\setlength{\emergencystretch}{3em}
\widowpenalty=10000
\clubpenalty=10000
\displaywidowpenalty=10000
\allowdisplaybreaks[1]
\setlist{nosep,leftmargin=1.5em}
\providecommand{\tightlist}{\setlength{\itemsep}{0pt}\setlength{\parskip}{0pt}}
\providecommand{\pandocbounded}[1]{#1}
\pagestyle{fancy}
\fancyhf{}
\fancyhead[L]{\sffamily\footnotesize\color{muted}OPEN PROBLEMS IN NUMERICAL LINEAR ALGEBRA}
\fancyhead[R]{\sffamily\bfseries\footnotesize\color{accent}SP-08}
\fancyfoot[L]{\parbox[b]{0.9\textwidth}{\sffamily\fontsize{8}{10}\selectfont\color{muted}Editorial ratings; a bounded search cannot certify that a problem remains unsolved.\\Literature check: 2026-09-08}}
\fancyfoot[R]{\sffamily\footnotesize\color{muted}\thepage}
\renewcommand{\headrulewidth}{0.3pt}
\renewcommand{\headrule}{\hbox to\headwidth{\color{accent}\leaders\hrule height \headrulewidth\hfill}}
\makeatletter
\renewcommand{\section}{\Needspace{5\baselineskip}\@startsection{section}{1}{0pt}{12pt plus 2pt minus 2pt}{4pt}{\sffamily\bfseries\large\color{ink}}}
\renewcommand{\subsection}{\Needspace{4\baselineskip}\@startsection{subsection}{2}{0pt}{9pt plus 1pt minus 1pt}{3pt}{\sffamily\bfseries\normalsize\color{ink}}}
\makeatother
\setcounter{secnumdepth}{0}
\begin{document}
{\sffamily\small\color{muted}Eigenvalues and inverse problems\par}
\vspace{3pt}
{\raggedright\hyphenpenalty=10000\exhyphenpenalty=10000\sffamily\bfseries\fontsize{21}{24}\selectfont\color{ink}Rank-two
maximizers of spectral spread on an entry interval\par}
\vspace{7pt}
{\sffamily\small\textbf{Difficulty:} challenging\quad\textbf{Importance:} interesting
to the community\par}
\vspace{4pt}
{\color{accent}\hrule height 0.8pt}
\vspace{6pt}
\textbf{Status:} open in general; several dimension and interval cases
proved.

\subsection{Problem statement}\label{problem-statement}

For an integer \(n\ge2\) and \(a\in[-1,1)\), let

\[
\mathcal S_n[a,1]=\{A\in\mathbb R^{n\times n}:A=A^T,\ a\le A_{ij}\le1
\text{ for all }i,j\}.
\]

The spread of \(A\) is \(s(A)=\lambda_{\max}(A)-\lambda_{\min}(A)\). The
Fallat--Xing conjecture asserts that there exists a matrix \(B\) of rank
exactly two, every entry of which belongs to \(\{a,1\}\), such that

\[
s(B)=\max_{A\in\mathcal S_n[a,1]}s(A).
\]

All diagonal entries are free to vary within the same interval as the
off-diagonal entries. The assertion concerns existence of a rank-two
maximizer; it does not assert that every maximizer has rank two. The
normalization \([a,1]\), \(-1\le a<1\), is the source's reduction of the
arbitrary nondegenerate real-interval problem. All dimensions and
intervals form one problem.

\subsection{Why it matters}\label{why-it-matters}

The conjecture reduces an eigenvalue optimization over a full symmetric
matrix family to a low-rank extremizer. It sharpens spectral bounds
obtainable from entrywise uncertainty alone.

\subsection{References}\label{references}

S. M. Fallat and Y. Xing,
\href{https://doi.org/10.1080/03081087.2012.703189}{\emph{On the spread
of certain normal matrices}}, Linear and Multilinear Algebra 60 (2012),
1391--1407, §2. N. J. Calkin, R. M. Corless, L. Gonzalez-Vega, J. R.
Sendra, and J. Sendra, \href{https://arxiv.org/abs/2510.15919}{\emph{On
the maximal spread of symmetric Bohemian matrices}}, 2025, §1 (displayed
Fallat--Xing conjecture), §§3.2, 7--8, and §10.

\subsection{Status check}\label{status-check}

Version 1 of the 2025 paper is the latest listed version checked. It
proves the conjecture for all \(a\in(-1,1)\) when \(2\le n\le7\), and
for \(a=0\) when \(2\le n\le8\) or \(3\) divides \(n\). The case
\(a=-1\) was already proved by Zhan. Its conclusion explicitly records a
gap in the authors' attempted general proof. The
\href{https://github.com/rcorless/BohemianSpread}{authors' supporting
repository} confirms the finite computational ranges. Searches for
\textit{Fallat Xing spread conjecture proof 2026} and later
maximal-spread papers found no general resolution. Restricting entries
to interval endpoints alone is already known and is not counted
separately.
\end{document}
